\chapter{Existencia y unicidad de soluciones}

\begin{definition}[Problema de valor inicial]
Sea $\displaystyle D \subset \R \times \R^{n} $ un abierto, $\displaystyle f : D \to \R^{n} $ continua y $\displaystyle \left(t_{0}, x_{0}\right) \in D $. Se llama \textbf{problema de valor inicial} al sistema
\[
\begin{cases}
x' = f\left(t,x\right) \\
x\left(t_{0}\right)=x_{0}
\end{cases}
.\]
\end{definition}

\begin{definition}[Solución]
	Una colección $\displaystyle x_{1}, \ldots, x_{n} : I\subset \R \to \R $ es una \textbf{solución} del sistema $\displaystyle x' = f\left(t, x\right) $ si $\displaystyle x_{i}\in \mathcal{C}^{1}\left(I, \R\right) $ y además $\displaystyle x_{i}' \left(t\right)=f_{i}\left(t, x_{1}\left(t\right), \ldots, x_{n}\left(t\right)\right) $, $\displaystyle \forall t \in I $, $\displaystyle \forall i \in \left\{ 1, \ldots, n\right\}  $. 
\end{definition}
Sea el PVI de la forma
\[
\begin{cases}
x' = f\left(t, x\right) \\
x\left(t_{0}\right)=x_{0}
\end{cases}
\]
con $\displaystyle f : D \subset \R \times \R^{n} \to \R^{n} $ y $\displaystyle \left(t_{0}, x_{0}\right)\in D $. Nos preguntamos si existe una solución $\displaystyle x \in \mathcal{C}^{1}\left(I \subset \R, \R^{n}\right) $ y si es única.
\begin{eg}
	Los siguientes ejemplos ilustran la cuestión.
	\begin{enumerate}
	\item Consideremos el PVI siguiente
\[
\begin{cases}
x' = a x \\
x\left(t_{0}\right)=x_{0}
\end{cases}
.\]
Tenemos que una solución es $\displaystyle x\left(t\right)=x_{0}e^{a\left(t-t_{0}\right)} $. Anteriormente vimos que esta solución es única. Por tanto, en este ejemplo tenemos existencia y unicidad.
\item Consideremos ahora el PVI
	\[
	\begin{cases}
	x' = 3x^{\frac{2}{3}} \\
	x\left(t_{0}\right)=x_{0}
	\end{cases}
	.\]
	Consideremos el caso en que $\displaystyle \left(t_{0}, x_{0}\right)=\left(0,0\right) $, entonces $\displaystyle x\left(t\right)=0 $ es solución. Por otro lado, por separación de variables obtenemos que también es solución la función $\displaystyle x\left(t\right)=t ^{3} $. Por otro lado, dados $\displaystyle a < 0 < b $, la función
	\[x\left(t\right) =
	\begin{cases}
		\left(x-a\right)^{3}, \; x \leq a \\
		0, \; a \leq x \leq b \\
		\left(x-b\right)^{3}, \; x \geq b
	\end{cases}
	,\]
	también es solución. 
	Por tanto, aunque hay existencia, no hay unicidad de soluciones. Si $\displaystyle t_{0}\neq 0 $ sí obtenemos unicidad de soluciones.
\item Consideremos ahora el PVI
	\[
	\begin{cases}
	x' = \frac{x}{t} \\
	x\left(t_{0}\right)=x_{0}
	\end{cases}
	.\]
Tenemos que podemos distinguir varios casos:
\begin{itemize}
\item Si $\displaystyle t_{0}\neq 0 $, entonces la solución única es $\displaystyle x\left(t\right)=\frac{x_{0}}{t_{0}}t $. 
\item Si $\displaystyle \left(t_{0}, x_{0}\right)=\left(0,0\right) $, tenemos que la solución es $\displaystyle x\left(t\right)=kt $, $\displaystyle \forall k \in \R $, por lo que la solución no es única. 
\item Si $\displaystyle t_{0}=0 $ y $\displaystyle x_{0}\neq 0 $, tenemos que no puede haber solución.
\end{itemize}
	\end{enumerate}
\end{eg}
\section{Formulación integral del PVI}
Por el Teorema Fundamental del Cálculo tenemos que 
\[
\begin{cases}
x'\left(t\right)=f\left(t, x\left(t\right)\right) \\
x\left(t_{0}\right)=x_{0}
\end{cases}
\Rightarrow x\left(t\right)=x_{0}+\int^{t}_{t_{0}} f\left(s, x\left(s\right)\right) \; d s.\]
\begin{prop}
Sea $\displaystyle f : D \subset \R \times \R^{n} \to \R^{n} $ continua y $\displaystyle \left(t_{0}, x_{0}\right)\in D $. Entonces son equivalentes
\begin{enumerate}
\item $x\in \mathcal{C}^{1}\left(I, \R^{n}\right) \; \text{con}\; 
\begin{cases}
x'\left(t\right)=f\left(t,x\left(t\right)\right) \\
x\left(t_{0}\right)=x_{0}. 
\end{cases}$.
\item $\displaystyle x \in \mathcal{C}\left(I, \R^{n}\right) \; \text{con}\; x\left(t\right)=x_{0} + \int^{t}_{t_{0}} f\left(s, x\left(s\right)\right) \; d s $.
\end{enumerate}
\end{prop}
\begin{proof}
Consideramos ambas implicaciones. 
\begin{description}
\item[($\Rightarrow$)] Claramente tenemos que $\displaystyle x \in \mathcal{C}\left(I, \R^{n}\right) $. Integrando tenemos que 
	\[x\left(t\right)-x_{0}=\int^{t}_{t_{0}} x'\left(s\right) \; d s = \int^{t}_{t_{0}} f\left(s, x\left(s\right)\right) \; d s .\]
\item[($\Leftarrow$)] Como $\displaystyle x \in \mathcal{C}\left(I, \R^{n}\right) $ tenemos que $\displaystyle f\left(t, x\left(t\right)\right) \in \mathcal{C}\left(I, \R^{n}\right)$. Por el TFC tenemos que 
	\[\int^{t}_{t_{0}} f\left(s, x\left(s\right)\right) \; d s \in \mathcal{C}^{1}\left(I, \R^{n}\right) .\]
	Así, de esta forma 
	\[x\left(t\right)=x_{0} + \int^{t}_{t_{0}} f\left(s, x\left(s\right)\right) \; d s \in \mathcal{C}^{1}\left(I, \R^{n}\right) .\]
Así, tenemos que $\displaystyle x \in \mathcal{C}^{1}\left(I, \R^{n}\right) $. Por otro lado, 
\[x\left(t_{0}\right)=x_{0} + \int^{t_{0}}_{t_{0}} f\left(s, x\left(s\right)\right) \; d s = x_{0} .\]
Finalmente, derivando $\displaystyle x\left(t\right) $ obtenemos que $\displaystyle x'\left(t\right)=f\left(t, x\left(t\right)\right) $ por el TFC. 
\end{description}
\end{proof}
\begin{notation}[Operador de Picard]
A partir de ahora vamos a definir el operador $\displaystyle T $ actuando en $\displaystyle \mathcal{C}\left(I, \R^{n}\right) $ como 
\[Ty\left(t\right):= x_{0} + \int^{t}_{t_{0}} f\left(s, y\left(s\right)\right) \; d s .\]
A este operador también se le llama \textbf{operador de Picard}. Por la proposición anterior, resolver el PVI es equivalente a encontrar $\displaystyle x \in \mathcal{C}\left(I, \R^{n}\right) $ tal que $\displaystyle Tx = x $. 
\end{notation}
\section{Teorema del punto fijo de Banach}
\begin{definition}[Espacio de Banach]
Diremos que $\displaystyle \left(E, \left| \cdot \right|\right) $ es un \textbf{espacio de Banach} si es un espacio vectorial normado y completo.
\end{definition}
\begin{eg}
Es sencillo ver que $\displaystyle \R^{n} $ es un espacio de banach con todas las normas de la forma
\[ \left|z\right|_{p}:=\left(\sum^{n}_{i = 1} \left|z_{i}\right|^{p}\right)^{\frac{1}{p}}, \; p \in [1,\infty) .\]
También lo es con la norma infinito.
\end{eg}
\begin{prop}
	Sea $\displaystyle I \subset \R $ un intervalo cerrado y acotado. Entonces $\displaystyle E = \mathcal{C}\left(I, \R^{n}\right) $ es un espacio de Banach con la norma $\displaystyle \|x\|_{\infty}= \max \left\{ \|x\left(t\right)\| \; : \; t \in I\right\}  $.
\end{prop}
\begin{proof}
	En primer lugar, es fácil comprobar que $\displaystyle E $ es un espacio vectorial sobre $\displaystyle \R $. Veamos ahora que $\displaystyle \| \cdot \|_{\infty} $ es una norma en $\displaystyle E $:
	\begin{itemize}
	\item La aplicación está bien definida puesto que al tratarse de funciones continuas definidas sobre un intervalo cerrado y acotado, siempre alcanzan su máximo.
	\item Es sencillo ver que $\displaystyle \| x\left(t\right)\|_{\infty} \geq 0 $, $\displaystyle \forall x\left(t\right) \in \mathcal{C}\left(I, \R^{n}\right) $. 
	\item Si $\displaystyle \|x\left(t\right)\|_{\infty}=0 $, tenemos que $\displaystyle \max \left\{ \left|x\left(t\right)\right| \; : \; t \in I\right\} =0  $, por lo que $\displaystyle \left|x\left(t\right)\right|\leq 0 $, $\displaystyle \forall t \in I $ y debe ser que $\displaystyle x \equiv 0 $.
	\item Dado $\displaystyle \lambda \in \R $ y $\displaystyle x\left(t\right)\in \mathcal{C}\left(I, \R^{n}\right) $, 
		\[
		\begin{split}
			\|\lambda x\left(t\right)\|_{\infty} = & \max \left\{ \left|\lambda x\left(t\right)\right| \; : \; t \in I\right\} = \max \left\{ \left|\lambda \right| \left|x\left(t\right)\right| \; : \; t \in I\right\} \\
			= &  \left|\lambda \right| \max \left\{ \left|x\left(t\right)\right|\; : \; t \in I\right\} =\left|\lambda \right|  \|x\left(t\right)\|_{\infty} .
		\end{split}
		\]
	\item Sean $\displaystyle x\left(t\right), y\left(t\right), z\left(t\right) \in \mathcal{C}\left(I, \R^{n}\right) $, entonces
		\[
		\begin{split}
			\|x\left(t\right)+z\left(t\right)\|_{\infty} = & \max \left\{ \left|x\left(t\right) + z\left(t\right)\right| \; : \; t \in I\right\} = \max \left\{ \left|x\left(t\right)-y\left(t\right)+y\left(t\right)-z\left(t\right)\right| \; : \; t \in I\right\}\\
			\leq & \max \left\{ \left|x\left(t\right)-y\left(t\right)\right|\; : \; t \in I\right\} + \max \left\{ \left|y\left(t\right)-z\left(t\right)\right| \; : \; t \in I\right\} \\
			= &  \|x\left(t\right)-y\left(t\right)\|_{\infty} + \|x\left(t\right)-z\left(t\right)\|_{\infty}.
		\end{split}
		\]
	\end{itemize}
	Así, tenemos que $\displaystyle \left(E, \| \cdot \|_{\infty}\right) $ es un espacio normado. Finalmente, veamos que es completo. Tenemos que una sucesión $\displaystyle \left\{ f_{m}\right\} _{m\in\N}\subset \mathcal{C}\left(I, \R^{n}\right) $ converge a $\displaystyle f $ si y solo si $\displaystyle \|f_{m}-f\|_{\infty} \to 0 $. Esto es, para cada $\displaystyle \epsilon > 0 $ debe existir un $\displaystyle m_{0} \in \N $ tal que $\displaystyle \forall m \geq m_{0} $ se tiene que 
	\[\|f_{m}\left(t\right)-f\left(t\right)\|_{\infty} < \epsilon .\]
	Sea $\displaystyle \left\{ f_{m}\right\} _{m\in\N} $ una sucesión de Cauchy y sea $\displaystyle t_{0} \in I $. Tenemos que la sucesión $\displaystyle \left\{ f_{m}\left(t_{0}\right)\right\} _{m\in\N}\subset \R^{m} $ es de Cauchy y como $\displaystyle \R^{n} $ es completo, es una sucesión convergente. Así, podemos definir $\displaystyle f : I \to \R^{n} $ tal que $\displaystyle \forall t \in I $, 
	\[f\left(t\right):=\lim_{m \to \infty}f_{m}\left(t\right) .\]
	Veamos que $\displaystyle \left\{ f_{m}\right\} _{m\in\N} $ converge a $\displaystyle f $ con la norma $\displaystyle \| \cdot \|_{\infty} $. Como $\displaystyle \left\{ f_{m}\right\} _{m\in\N} $ es de Cauchy, dado $\displaystyle \epsilon > 0 $ existe $\displaystyle m_{0} \in \N $ tal que $\displaystyle \forall m,l \geq m_{0} $ se tiene que $\displaystyle \|f_{m}-f_{l}\|_{\infty} < \epsilon  $, por lo que 
	\[ \|f_{m}\left(t\right)-f_{l}\left(t\right)\| < \epsilon, \; \forall t \in I .\]
	Tomando el límite cuando $\displaystyle l \to \infty $ obtenemos que
	\[ \|f_{m}\left(t\right)-f\left(t\right)\|<\epsilon, \; \forall t \in I ,\]
	por lo que $\displaystyle \|f_{m}-f\|_{\infty}<\epsilon  $ y la sucesión $\displaystyle \left\{ f_{m}\right\} _{m\in\N} $ converge a $\displaystyle f $ en $\displaystyle \| \cdot \|_{\infty} $. Faltaría ver que la función $\displaystyle f $ es continua, pero esto es muy sencillo y se deja como ejercicio para el lector.
\end{proof}
\begin{definition}[Aplicación contractiva]
Sea $\displaystyle E $ un espacio de Banach. Diremos que $\displaystyle f: X \subset E \to X $ es una \textbf{aplicación contractiva} si existe $\displaystyle K \in [0,1) $ tal que 
\[ \|f\left(x\right)-f\left(y\right)\|\leq K \|x-y\|, \; \forall x,y \in X .\]
\end{definition}
\begin{observation}
Si $\displaystyle f $ es contractiva, entonces es uniformemente continua.
\end{observation}
\begin{theorem}[Teorema del punto Banach]
Sea $\displaystyle \emptyset \neq X \subset E$ cerrado, $\displaystyle E $ un espacio de Banach y $\displaystyle T : X \to X $ una aplicación contractiva. Entonces, existe un único punto fijo $\displaystyle x \in X $ y se puede encontrar como el límite de la sucesión $\displaystyle x_{j+1}=Tx_{j} $, $\displaystyle \forall x_{0} \in X $. 
\end{theorem}
\begin{proof}
Demostramos primero la unicidad. Supongamos que existen dos puntos fijos $\displaystyle x,y \in X $. Entonces, existe $\displaystyle K \in [0,1) $ tal que 
\[ \|x - y \| = \|Tx - Ty\| \leq K \|x-y\| .\]
Si $\displaystyle K = 0 $ hemos terminado porque tenemos que $\displaystyle x = y $. Si $\displaystyle K \neq 0 $ tenemos que $\displaystyle 0 < K < 1 $, por lo que
\[\|x-y\| \leq K\|x-y\| < \|x-y\| ,\]
esto es una contradicción si $\displaystyle x \neq y $. \\
Ahora demostramos la existencia. Dado $\displaystyle x_{0} \in X $, definimos la sucesión $\displaystyle x_{j+1}=Tx_{j} $. Veamos que se trata de una sucesión de Cauchy. Dados $\displaystyle i > j $ tenemos que 
\[
\begin{split}
	\|x_{i}-x_{j}\| = & \|x_{i}-x_{i-1} + x_{i-1} + \cdots -x_{j}\| \leq \|x_{i}-x_{i-1}\|+ \cdots + \|x_{j+1}-x_{j}\| \\
	= & \|Tx_{i-1}-Tx_{i-2}\| + \|x_{i-1}-x_{i-2}\| + \cdots + \|x_{j+1}-x_{j}\| \\
	\leq & \left(K^{i-1} + K^{i-2} + \cdots + K^{j}\right)\|x_{1}-x_{0}\| \leq \left(\sum^{\infty}_{k = j}K^{j}\right)\|x_{1}-x_{0}\| \\
	= & \frac{K^{j}}{1-K} \|x_{1}-x_{0}\| \xrightarrow{j \to \infty} 0.
\end{split}
\]
Así, como $\displaystyle \left\{ x_{j}\right\} _{j \in \N} $ es una sucesión de Cauchy y $\displaystyle E $ es completo, debe ser que existe $\displaystyle x \in E $ tal que $\displaystyle x_{j}\to x $. Además, como $\displaystyle X $ es cerrado debe ser que $\displaystyle x \in X $. Por otro lado, como $\displaystyle T $ es contractiva tenemos que es continua y por tanto podemos tomar el límite a ambos lados de la igualdad $\displaystyle x_{j+1} = Tx_{j} $, obteniendo $\displaystyle x = Tx $. 
\end{proof}

\begin{colorary}
Es suficiente que una potencia \footnote{Entendemos por potencia componer $\displaystyle T $, $\displaystyle n $ veces consigo mismo para algún $\displaystyle n\in \N $.} de $\displaystyle T $ sea contractiva para que se verifique el resultado.
\end{colorary}
\begin{proof}
Supongamos que existe $m \in \mathbb{N}$ tal que $T^m$ es contractiva. Es decir, existe
$K \in [0,1)$ tal que
\[\|T^m x - T^m y\| \leq K \|x-y\|,\; \forall x,y \in X.\]
Definimos $\displaystyle S := T^m $. Como $X$ es cerrado en un espacio de Banach, $X$ es completo. Por el teorema del punto fijo de Banach, $S$ tiene un único punto fijo $x^\ast \in X$, es decir, $\displaystyle T^m x^\ast = x^\ast $.
Veamos que $x^\ast$ es también punto fijo de $T$. En efecto,
\[T^m(Tx^\ast)=T(T^m x^\ast)=Tx^\ast.\]
Por tanto, $Tx^\ast$ es un punto fijo de $T^m$. Como $T^m$ tiene un único punto fijo, se sigue que $\displaystyle Tx^\ast = x^\ast $. 
Luego $x^\ast$ es un punto fijo de $T$.
Además, este punto fijo es único. En efecto, si $y \in X$ satisface $\displaystyle Ty = y$, entonces
\[T^m y=y.\]
Por tanto, $y$ es también un punto fijo de $T^m$. Por unicidad del punto fijo de $T^m$ se tiene que $\displaystyle y=x^\ast $. 
Finalmente, veamos que para todo $x_0 \in X$ la sucesión definida por $\displaystyle x_{j+1}=Tx_j $
converge a $x^\ast$. Como $S=T^m$ es contractiva, por el teorema de Banach tenemos que
\[S^k z \longrightarrow x^\ast, \; \forall z \in X.\]
Dado $j \in \mathbb{N}$, por la división euclídea existen $\displaystyle k,r \in \N $ tales que
\[j=km+r,\quad 0 \leq r < m.\]
Entonces
\[x_j=T^j x_0=T^{km+r}x_0=(T^m)^k(T^r x_0)=S^k(T^r x_0).\]
Fijado $r \in \{0,\dots,m-1\}$, como $T^r x_0 \in X$, se tiene
\[S^k(T^r x_0)\longrightarrow x^\ast.\]
Por tanto,
\[x_{km+r}\longrightarrow x^\ast\quad\forall r=0,\dots,m-1.\]
Así, todas las subsucesiones correspondientes a los distintos restos módulo $m$
convergen al mismo límite $x^\ast$. Como sólo hay un número finito de ellas,
la sucesión completa converge:
\[
x_j=T^j x_0 \longrightarrow x^\ast.
\]
\end{proof}
\section{Condiciones para el teorema de existencia y unicidad}
\begin{lema}
	Sea $\displaystyle f \in \mathcal{C}\left(\left[a,b\right] , \R^{n}\right) $. Entonces,
	\[ \left\|\int^{b}_{a} f\left(t\right) \; dt\right\| \leq \int^{b}_{a} \|f\left(t\right)\| \; dt.\]
\end{lema}
\begin{proof}
Como $\displaystyle f\left(t\right)=\left(f_{1}\left(t\right), \ldots, f_{n}\left(t\right)\right)\in \R^{n} $, la integral de $\displaystyle f $ se define como la integral componente a componente. Es decir, 
\[\int^{b}_{a} f\left(t\right) \; dt : = \left(\int^{b}_{a} f_{1}\left(t\right) \; dt, \ldots, \int^{b}_{a} f_{n}\left(t\right) \; dt\right) .\]
Así, tendremos que 
\[
\begin{split}
	\left\|\int^{b}_{a} f\left(t\right) \; dt\right\|^{2} = & \sum^{n}_{i = 1}\left(\int^{b}_{a} f_{i}\left(t\right) \; dt\right)^{2} = \sum^{n}_{i = 1}\left(\int^{b}_{a} f_{i}\left(t\right) \; dt\right)\left(\int^{b}_{a} f_{i}\left(s\right) \; d s\right)\\
	= & \sum^{n}_{i = 1}\int \int _{\left[a,b\right]\times\left[a,b\right]}f_{i}\left(t\right)f_{i}\left(s\right) \; dt  \; d s = \int \int _{\left[a,b\right]\times\left[a,b\right]}\left(\sum^{n}_{i = 1}f_{i}\left(t\right)f_{i}\left( s\right)\right) \; dt \; d s,
\end{split}
\]
aplicando la desigualdad Cauchy-Swarz,
\[
	\leq  \int \int _{\left[a,b\right]\times\left[a,b\right]} \|f\left(t\right)\|\|f\left(s\right)\| \; dt \; d s =\left(\int^{b}_{a} \|f\left(t\right)\| \; dt\right)\left(\int^{b}_{a} \|f\left(s\right)\| \; d s\right)=\left(\int^{b}_{a} \|f\left(t\right)\| \; dt\right)^{2}.
\]
\end{proof}

Para que la aplicación 
\[Tx\left(t\right)=x_{0}+\int^{t}_{t_{0}} f\left(s, x\left(s\right)\right) \; d s ,\]
sea contractiva, necesitamos que existe $\displaystyle K \in [0,1) $ tal que 
\[\|Ty-Tz\|_{\infty} \leq K \|y-z\|_{\infty} .\]
Veamos qué tiene que suceder con $\displaystyle f $ para que $\displaystyle T $ sea contractiva:
\[
\begin{split}
	\|Ty-Tz\|_{\infty} = & \max _{t \in I} \|Ty\left(t\right)-Tz\left(t\right)\| = \max_{t \in I}\left\|\int^{t}_{t_{0}} f\left(s,y\left(s\right)\right)-f\left(s, z\left(s\right)\right) \; d s\right\|  \\
	\leq & \max _{t \in I}\int^{t}_{t_{0}} \|f\left(s,y\left(s\right)\right)-f\left(s, z\left(s\right)\right) \; d s.
\end{split}
\]
Tenemos que poder comparar esto con 
\[K \|y -z\|_{\infty} = K\max_{t \in I}\|y\left(t\right)-z\left(t\right)\| .\]
Tendríamos que poder comparar entonces $\displaystyle \|f\left(s, y\left(s\right)\right)-f\left(s, z\left(s\right)\right)\| $ y $\displaystyle \|y\left(s\right)-z\left(s\right)\| $.
\begin{definition}[Función Lipschitz]
Sea $\displaystyle f : \Omega \subset \R^{n} \to \R^{m} $.
\begin{itemize}
\item Diremos que $\displaystyle f $ es \textbf{globalmente Lipschitz} si existe $\displaystyle L \geq 0 $ tal que 
	\[ \|f\left(y\right)-f\left(z\right)\| \leq L \|y-z\|, \; \forall y,z \in \Omega  .\]
\item Diremos que es \textbf{localmente Lipschitz} si $\displaystyle \forall x \in \Omega  $, existe un entorno $\displaystyle V_{x}\subset \Omega $ y $\displaystyle L _{x}\geq 0 $ tales que 
	\[\|f\left(y\right)-f\left(z\right)\| \leq L_{x}\|y-z\|, \; \forall y,z \in V_{x} .\]
\end{itemize}
\end{definition}
\begin{prop}
Sea $\displaystyle f: \Omega\subset \R^{n} \to \R^{m} $ localmente Lipschitz y $\displaystyle K \subset \Omega  $ compacto. Entonces $\displaystyle f $ es globalmente Lipschitz en $\displaystyle K $. 
\end{prop}
\begin{proof}
Supongamos que no es cierto. Entonces para cada $\displaystyle k \in \N $ tenemos que existen $\displaystyle x_{k}, y_{k}\in K $ distintos tales que 
\[ \|f\left(x_{k}\right)-f\left(y_{k}\right)\| > k \|x_{k}-y_{k}\| .\]
Como $\displaystyle K $ es compacto, existen subsucesiones $\displaystyle \left\{ x_{k_{j}}\right\} _{j \in \N} $ e $\displaystyle \left\{ y_{k_{j}}\right\} _{j \in \N} $ convergentes, por lo que existen $\displaystyle \overline{x}, \overline{y}\in K $ con $\displaystyle x_{k_{j}}\to \overline{x} $ e $\displaystyle y_{k_{j}}\to \overline{y} $. 
Como $\displaystyle f $ es continua y $\displaystyle K $ compacto, entonces $\displaystyle f $ está acotada en $\displaystyle K $, de forma que 
\[\|x_{k_{j}}-y_{k_{j}}\|< \frac{\|f\left(x_{k_{j}}\right)-f\left(y_{k_{j}}\right)\|}{k_{j}}<\frac{2\max_{K} \|f\|}{k_{j}} .\]
Tomando límites obtenemos que 
\[\|\overline{x}-\overline{y}\| \leq \lim_{j \to \infty}\frac{2\max_{K} \|f\|}{k_{j}} = 0 ,\]
por lo que debe ser que $\displaystyle \overline{x}=\overline{y} $. Como $\displaystyle f $ es localmente Lipschitz, existe $\displaystyle V_{\overline{x}}\subset \Omega  $ entorno de $\displaystyle \overline{x} $ y $\displaystyle L_{\overline{x}} \geq 0$ tal que
\[\|f\left(x\right)-f\left(y\right)\|\leq L_{\overline{x}}\|x-y\|, \; \forall x,y\in V_{\overline{x}} .\]
Como las sucesiones $\displaystyle \left\{ x_{k_{j}}\right\} _{j \in \N} $ e $\displaystyle \left\{ y_{k_{j}}\right\} _{j \in \N} $ convergen al mismo límite, existe $\displaystyle N \in \N $ tal que $\displaystyle x_{k_{j}}, y_{k_{j}} \in V_{\overline{x}} $, $\displaystyle \forall j \geq N $. 
Por tanto,
\[k_{j}\|x_{k_{j}}-y_{k_{j}}\| < \|f\left(x_{k_{j}}\right)-f\left(y_{k_{j}}\right)\| \leq L_{\overline{x}}\|x_{k_{j}}-y_{k_{j}}\|, \; \forall j \geq N .\]
Como $\displaystyle x_{k_{j}}\neq y_{k_{j}} $ se concluye que $\displaystyle L_{\overline{x}}>k_{j} $, $\displaystyle \forall j \geq N $, que es una contradicción. 
\end{proof}

\begin{prop}
Sea $\displaystyle I = \left(a,b\right) $ y $\displaystyle f \in \mathcal{C}^{1}\left(I, \R\right) $. Entonces $\displaystyle f $ es Lipschitz en $\displaystyle I $ si y solo si $\displaystyle f' $ está acotada en $\displaystyle I $. 
\end{prop}
\begin{proof}
Consideramos ambas implicaciones.
\begin{description}
\item[($\displaystyle \Rightarrow $)] Cogemos $\displaystyle t \in I $ y $\displaystyle h \in \R $ suficientemente pequeño tal que $\displaystyle t + h \in I $. Como $\displaystyle f $ es Lipschitz, existe $\displaystyle L \geq0 $ tal que
	\[ \left|f\left(t + h\right)-f\left(t\right)\right| \leq L \left|h\right| \Rightarrow \left|\frac{f\left(t + h\right)-f\left(t\right)}{h}\right|\leq L.\]
Así, nos queda que si tomamos el límite cuando $\displaystyle h \to 0 $, $\displaystyle \left|f'\left(t\right)\right|\leq L $. 
\item[($\displaystyle \Leftarrow $)] Tomamos $\displaystyle x,y \in I $ con $\displaystyle x > y $, entonces
	\[f\left(x\right)-f\left(y\right)=\int^{x}_{y} f'\left(t\right) \; dt .\]
	Tomando valores absolutos,
	\[ \left|f\left(x\right)-f\left(y\right)\right| = \left|\int^{x}_{y} f'\left(t\right) \; dt\right|\leq \int^{x}_{y} \left|f'\left(t\right)\right| \; dt \leq \sup_{t \in I} \left|f'\left(t\right)\right| \left|x-y\right| .\]
	Esta implicación también se puede demostrar con el teorema del valor medio.
\end{description}
\end{proof}

\begin{observation}
Existen funciones Lipschitz que no son $\displaystyle \mathcal{C}^{1} $. Un ejemplo es la función $\displaystyle f\left(x\right) = \left|x\right| $. 
\end{observation}

\begin{eg}[Función uniformemente continua que no es Lipschitz]
Consideremos $\displaystyle f : \left(0,1\right)\to \R : t \to \sqrt{t} $. Tenemos que $\displaystyle f'\left(t\right)=\frac{1}{2\sqrt{t}} $ no está acotada en $\displaystyle \left(0,1\right) $. Por el teorema anterior tenemos que no es Lipschitz, pero veamos que no lo es a partir de la definición. Para ello, supongamos que lo es, es decir, existe $\displaystyle L \geq 0$ tal que 
\[ \left|\sqrt{x}-\sqrt{y}\right|\leq L \left|x-y\right| .\]
Si tomamos las sucesiones $\displaystyle x_{k}=\frac{1}{k^{2}} $ e $\displaystyle y_{k}=\frac{1}{\left(k + 1\right)^{2}} $, tenemos que 
\[ \left|\sqrt{x_{k}}-\sqrt{y_{k}}\right| = \left|\frac{1}{k}-\frac{1}{k + 1}\right|=\frac{1}{k\left(k+1\right)}\leq L\frac{2k+1}{k^{2}\left(k+1\right)^{2}} .\]
Tendremos entonces que 
\[L \geq \frac{k\left(k + 1\right)}{2k+1}, \; \forall k \in \N .\]
Sin embargo, esta sucesión tiende a $\displaystyle \infty $, por lo que hemos llegado a una contradicción. 
\end{eg}
\begin{prop}
	Sea $\displaystyle f\in \mathcal{C}^{1}\left(\Omega\subset \R^{n}, \R^{m}\right) $ con $\displaystyle \Omega  $ convexo. Entonces, $\displaystyle f $ es Lipschitz en $\displaystyle \Omega  $ si y solo si existe $\displaystyle C \geq 0 $ tal que $\displaystyle \left|\frac{\partial f_{k}}{\partial x_{j}}\left(x\right)\right| \leq C $, $\displaystyle \forall k \in \left\{ 1, \ldots, m\right\}  $, $\displaystyle \forall j \in \left\{ 1, \ldots, n\right\}  $, $\displaystyle \forall x \in \Omega  $. 
\end{prop}
\begin{proof}
	Podemos escribir $\displaystyle f = \left(f_{1}, \ldots, f_{m}\right) $. 
	\begin{description}
	\item[($\displaystyle \Rightarrow $)] Supongamos que existe $\displaystyle L \geq 0 $ tal que 
		\[\|f\left(x\right)-f\left(y\right)\|\leq L\|x-y\|, \; \forall x,y\in \Omega  .\]
		Fijemos $\displaystyle x \in \Omega  $ y $\displaystyle j\in \left\{ 1, \ldots, n\right\}  $. Como $\displaystyle \Omega  $ es abierto, existe $\displaystyle h\in \R/ \left\{ 0\right\}  $ suficientemente pequeño tal que $\displaystyle x + he_{j}\in \Omega  $. Tendremos que 
		\[\|f\left(x+he_{j}\right)-f\left(x\right)\|\leq L\|he_{j}\|=L \left|h\right| \Rightarrow \left\|\frac{f\left(x + he_{j}\right)-f\left(x\right)}{h}\right\|\leq L.\]
		Así, tomando el límite cuando $\displaystyle h \to 0 $, tenemos que $\displaystyle \left\|\frac{\partial f}{\partial x_{j}}\right\| \leq L$. 
		Como 
		\[\frac{\partial f}{\partial x_{j}}=\left(\frac{\partial f_{1}}{\partial x_{j}}, \ldots, \frac{\partial f_{m}}{\partial x_{j}}\right) ,\]
		tenemos que cada componente satisface
		\[ \left|\frac{\partial f_{k}}{\partial x_{j}}\right|\leq \left\|\frac{\partial f}{\partial x_{j}}\right\|\leq L .\]
	\item[($\displaystyle \Leftarrow $)] Supongamos ahora que existe $\displaystyle C \geq 0 $ tal que $\displaystyle \left|\frac{\partial f_{k}}{\partial x_{j}}\right|\leq C $, $\displaystyle \forall x \in \Omega  $, $\displaystyle \forall j \in \left\{ 1, \ldots, n\right\}  $, $\displaystyle \forall k \in \left\{ 1, \ldots, m\right\}  $. 
		Si $\displaystyle x,y\in \Omega  $, por ser $\displaystyle \Omega  $ convexo tenemos que $\displaystyle \left[x,y\right] \subset \Omega  $. Es decir, $\displaystyle x + t\left(y-x\right)\in \Omega  $, $\displaystyle \forall t \in \left[0,1\right]  $. Así, para cada componente $\displaystyle f_{k} $ definimos
		\[g_{k}\left(t\right)=f_{k}\left(x+t\left(y-x\right)\right) .\]
		Es sencillo ver que $\displaystyle g_{k}\left(1\right)-g_{k}\left(0\right)=f_{k}\left(y\right)-f_{k}\left(x\right) $. Aplicando la regla de la cadena,
		\[g'_{k}\left(t\right)=\sum^{n}_{j = 1}\frac{\partial f_{k}}{\partial x_{j}}\left(x+t\left(y-x\right)\right)\left(y_{j}-x_{j}\right) .\]
		Por el teorema fundamental del cálculo,
		\[f_{k}\left(y\right)-f_{k}\left(x\right)=\int^{1}_{0} g'_{k}\left(t\right) \; dt .\]
	Por tanto, 
	\[ \left|f_{k}\left(y\right)-f_{k}\left(x\right)\right|\leq \int^{1}_{0} \sum^{n}_{j = 1} \left|\frac{\partial f_{k}}{\partial x_{j}}\left(x + t\left(y-x\right)\right)\right| \left|y_{j}-x_{j}\right| \; dt .\]
	Usando la hipótesis, tenemos que
	\[ \left|f_{k}\left(y\right)-f_{k}\left(x\right)\right|\leq C\sum^{n}_{j=1} \left|y_{j}-x_{j}\right| ,\]
	y por Cauchy-Schwarz,
	\[\sum^{n}_{j = 1} \left|y_{j}-x_{j}\right|\leq \sqrt{n} \|y-x\| .\]
	Por tanto, nos queda que $\displaystyle \left|f_{k}\left(y\right)-f_{k}\left(x\right)\right|\leq C\sqrt{n} \|y-x\| $. Ahora, para $\displaystyle f = \left(f_{1}, \ldots, f_{m}\right) $ tendremos que 
	\[\|f\left(y\right)-f\left(x\right)\|^{2}=\sum^{m}_{k = 1} \left|f_{k}\left(y\right)-f_{k}\left(x\right)\right|^{2}\leq \sum^{m}_{k = 1} C^{2}n \|y-x\|^{2}=mnC^{2}\|y-x\|^{2} .\]
	Tomando raíces obtenemos que $\displaystyle \|f\left(y\right)-f\left(x\right)\|\leq C\sqrt{nm}\|y-x\| $, por lo que $\displaystyle f $ es Lipschitz. 
	\end{description}
\end{proof}
\begin{observation}
La segunda implicación de la proposición anterior también se puede demostrar utilizando el teorema del valor medio en vez del teorema fundamental del cálculo.
\end{observation}

\begin{definition}[Lipschitz respecto a la segunda variable]
Dada $\displaystyle f : D\subset\R\times\R^{n}\to\R^{m} $, diremos que es
\begin{itemize}
\item \textbf{globalmente Lipschitz respecto a la segunda variable} si existe $\displaystyle L \geq0 $ 
	\[ \left\|f\left(t, x\right)-f\left(t,y\right)\right\| \leq L \left\|x-y\right\|, \; \forall \left(t,x\right), \left(t,y\right)\in D .\]
	
\item \textbf{localmente Lipschitz respecto a la segunda variable} si para $\displaystyle \left(t, x\right)\in D $ existe un entorno $\displaystyle U_{\left(t,x\right)} $ y $\displaystyle L_{\left(t,x\right)}\geq 0  $ tales que 
	\[ \left\|f\left(t,x\right)-f\left(t,y\right)\right\|\leq L_{\left(t,x\right)} \left\|x-y\right\| , \; \forall\left(t,x\right),\left(t,y\right)\in U_{\left(t,x\right)}.\]
\end{itemize}
\end{definition}
\begin{prop}
Dado $\displaystyle D \subset \R\times\R^{n} $ y $\displaystyle f \in \mathcal{C}\left(D, \R^{m}\right) $ tal que $\displaystyle \frac{\partial f_{k}}{\partial x_{j}} $ existen y están acotadas en $\displaystyle D $. Entonces $\displaystyle f $ es Lipschitz en $\displaystyle D $ respecto a la segunda variable. 
\end{prop}
\begin{proof}
	Podemos escribir $\displaystyle f = \left(f_{1}, \ldots, f_{m}\right) $ y tenemos que existe $\displaystyle C > 0 $ tal que $\displaystyle \left|\frac{\partial f_{k}}{\partial x_{j}}\left(t, x\right)\right|\leq C $, $\displaystyle \forall\left(t,x\right)\in D $, $\displaystyle \forall k \in \left\{ 1, \ldots, m\right\}  $, $\displaystyle \forall j \in \left\{ 1, \ldots, n\right\}  $. Fijemos $\displaystyle t $ y dos puntos $\displaystyle x = \left(x_{1}, \ldots, x_{n}\right) $ e $\displaystyle y = \left(y_{1}, \ldots, y_{n}\right) $. 
	Para pasar de $\displaystyle x $ a $\displaystyle y $ vamos cambiando una coordenada cada vez, es decir
	\[
	\begin{split}
		p_{0}=& \left(x_{1}, \ldots, x_{n}\right) \\
		p_{1} = & \left(y_{1}, x_{2}, \ldots, x_{n}\right)\\
		\vdots & \\
		p_{n} =& \left(y_{1}, \ldots, y_{n}\right).
	\end{split}
	\]
	Así, para cada componente $\displaystyle f_{k} $ tendremos que
	\[f_{k}\left(t,y\right)-f_{k}\left(t,x\right)=\sum^{n}_{j= 1}\left(f_{k}\left(t, p_{j}\right)-f_{k}\left(t,p_{j-1}\right)\right) .\]
	En el término $\displaystyle j $-ésimo sólo cambia la coordenada $\displaystyle x_{j} $. Aplicando el teorema del valor medio a la función de una variable correspondiente tenemos que existe $\displaystyle \xi_{j} $ tal que 
	\[f_{k}\left(t, p_{j}\right)-f_{k}\left(t, p_{j-1}\right)=\frac{\partial f_{k}}{\partial x_{j}}\left(t, \xi_{j}\right)\left(y_{j}-x_{j}\right) .\]
	De esta forma, tendremos que $\displaystyle \left|f_{k}\left(t, p_{j}\right)-f_{k}\left(t, p_{j-1}\right)\right|\leq C \left|y_{j}-x_{j}\right| $. Sumando,
	\[ \left|f_{k}\left(t,y\right)-f_{k}\left(t,x\right)\right|\leq C\sum^{n}_{j=1} \left|y_{j}-x_{j}\right| .\]
	Por Cauchy-Schwarz,
	\[\sum^{n}_{j=1} \left|y_{j}-x_{j}\right| \leq \sqrt{n} \|y-x\|.\]
	De esta fomra, $\displaystyle \left|f_{k}\left(t,y\right)-f_{k}\left(t,x\right)\right|\leq C\sqrt{n}\|y-x\| $. Si sumamos las $\displaystyle m $ componentes,
	\[\|f\left(t,y\right)-f\left(t,x\right)\|^{2} = \sum^{m}_{k=1} \left|f_{k}\left(t,y\right)-f_{k}\left(t,x\right)\right|^{2}\leq mC^{2}n\|y-x\|^{2} .\]
\[\Rightarrow \|f\left(t,y\right)-f\left(t,x\right)\|\leq C\sqrt{mn}\|y-x\| .\]	
\end{proof}
\begin{prop}
Si $\displaystyle f : D \subset \R \times \R^{n} \to \R^{m} $ es continua y localmente Lipschitz respecto a la segunda variable en $\displaystyle D $ abierto, entonces es globalmente Lipschitz respecto a la segunda variable en todo compacto $\displaystyle K \subset D $.
\end{prop}
\begin{proof}
	Para cada $\displaystyle \left(t_{0}, x_{0}\right)\in D $, existe un entorno $\displaystyle U_{\left(t_{0}, x_{0}\right)}\subset D $ y una constante $\displaystyle L_{\left(t_{0}, x_{0}\right)} \geq 0$ tales que
	\[\|f\left(t,x\right)-f\left(t,y\right)\| \leq L_{\left(t_{0}, x_{0}\right)}\|x-y\| , \; \forall\left(t,x\right), \left(t,y\right)\in U_{\left(t_{0}, x_{0}\right)}.\]
	Sea $\displaystyle K \subset D $ un compacto. Tenemos que la familia de entornos $\displaystyle \left\{ U_{\left(t_{0}, k_{0}\right)}\; : \; \left(t_{0}, x_{0}\right)\in K\right\}  $ es un recubrimiento abierto de $\displaystyle K $. Por ser $\displaystyle K $ compacto, podemos encontrar un subrecubrimiento finito de forma que 
	\[K \subset U_{1} \cup \cdots \cup U_{N} .\]
	Si $\displaystyle L_{i} $ es la constante de Lipschitz asociada a $\displaystyle U_{i} $ para $\displaystyle i \in \left\{ 1, \ldots, N\right\}  $, cogemos 
	\[L = \max \left\{ L_{1}, \ldots, L_{N}\right\}  .\]
	Así, si $\displaystyle \left(t,x\right) $ y $\displaystyle \left(t,y\right) $ están en el mismo $\displaystyle U_{i} $ tendremos que 
	\[\|f\left(t,x\right)-f\left(t,y\right)\|\leq L\|x-y\| .\]
	Sea $\displaystyle \delta > 0 $ tal que cualquier subconjunto de $\displaystyle K $ de diámetro menor que $\displaystyle \delta  $ está contenido en alguno de los $\displaystyle U_{i} $, con $\displaystyle i \in \left\{ 1, \ldots, N\right\}  $. Distinguimos dos casos:
	\begin{itemize}
	\item Si $\displaystyle \left(t,x\right), \left(t,y\right) \in K $ con $\displaystyle \|x-y\|<\delta  $, tendremos que 
		\[\|\left(t, x\right)-\left(t,y\right)\| = \|x-y\|<\delta  ,\]
		por lo que pertenecen al mismo $\displaystyle U_{i} $ y por lo que acabamos de ver tendremos que 
		\[\|f\left(t,x\right)-f\left(t,y\right)\|\leq L\|x-y\| .\]
	\item Si $\displaystyle \|x-y\|\geq \delta  $, como $\displaystyle f $ es continua y $\displaystyle K $ es compacto tenemos que existe $\displaystyle M \geq 0 $ tal que $\displaystyle \|f\left(t,x\right)\|\leq M $, $\displaystyle \forall\left(t,x\right)\in K $. Por la desigualdad triangular, 
		\[\|f\left(t,x\right)-f\left(t,y\right)\|\leq \|f\left(t,x\right)\|+\|f\left(t,y\right)\|\leq 2M .\]
		Como $\displaystyle \|x-y\|\geq \delta  $, tenemos que $\displaystyle 2M \leq \frac{2M}{\delta }\|x-y\| $. De esta forma, 
		\[\|f\left(t,x\right)-f\left(t,y\right)\|\leq \frac{2M}{\delta }\|x-y \|.\]
	\end{itemize}
	Así, basta con coger como constante de Lipschitz $\displaystyle L_{K}=\max \left\{ L, \frac{2M}{\delta }\right\}  $ y tendremos que 
	\[\|f\left(t,x\right)-f\left(t,y\right)\|\leq L_{K}\|x-y\|, \; \forall\left(t,x\right), \left(t,y\right)\in K .\]
\end{proof}
\section{Teorema de existencia y unicidad}

\begin{eg}[Funcionamiento del argumento]
	Consideremos el siguiente PVI:
	\[
	\begin{cases}
	x' = 2tx \\
	x\left(0\right)=1
	\end{cases}
	.\]
Consideremos el operador $\displaystyle T $ definido anteriormente,
\[\left(Ty\right)\left(t\right)= 1 + \int^{t}_{0} 2sy\left(s\right) \; d s .\]
Aplicamos el teorema del punto fijo para encontrar un punto fijo de $\displaystyle T $. Como $\displaystyle y_{0} $ cogemos la función nula. Así, tendremos que 
\[y_{1}\left(t\right)=\left(Ty_{0}\right)\left(t\right)=1 + \int^{t}_{0} 2s y_{0}\left(s\right) \; d s = 1 + 0 = 1 .\]
\[y_{2}\left(t\right)=\left(Ty_{1}\right)\left(t\right)=1 + \int^{t}_{0} 2sy_{1}\left(s\right) \; d s = 1 + t ^{2}.\]
\[y_{3}\left(t\right)=\left(Ty_{2}\right)\left(t\right)=1 + \int^{t}_{0} 2s\left(1 + s^{2}\right) \; d s = 1 + t ^{2} + \frac{t ^{4}}{2} .\]
En general, 
\[y_{k}\left(t\right)=\sum^{k-1}_{j = 0}\frac{t ^{2j}}{j!} \Rightarrow y\left(t\right)=\sum^{\infty}_{j = 0}\frac{t ^{2j}}{j!}=e^{t ^{2}} .\]
\end{eg}

\begin{theorem}[Teorema de Picard-Lindelöf local]
Sea $\displaystyle f : D \to \R^{n} $ continua y localmente Lipschitz respecto a la segunda variable en un abierto $\displaystyle D \subset \R \times \R^{n} $. Dado $\displaystyle \left(t_{0}, x_{0}\right)\in D $ existe $\displaystyle \epsilon > 0 $ tal que el PVI
\[
\begin{cases}
x' = f\left(t,x\right) \\
x\left(t_{0}\right)=x_{0}
\end{cases}
,\]
admite una única solución definida en el intervalo $\displaystyle I_{\epsilon } =\left[t_{0}-\epsilon , t_{0} + \epsilon \right]  $. 
\end{theorem}
\begin{proof}
Vamos a aplicar el teorema del punto fijo de Banach de forma adecuada al operador $\displaystyle T $ de la formulación integral del PVI. Si $\displaystyle f \equiv 0 $ el resultado es trival. Así, supongamos que no lo es. Para aplicar el TPF necesitamos encontrar conjuntos $\displaystyle E $ y $\displaystyle X \subset E $ adecuados tales que $\displaystyle T : X \subset E \to X$. Recordamos que $\displaystyle T $ viene dada por
\[\left(Ty\right)\left(t\right)=x_{0} + \int^{t}_{t_{0}} f\left(s, y\left(s\right)\right) \; d s .\]
Como $\displaystyle D $ es abierto, dado $\displaystyle \left(t_{0}, x_{0}\right)\in D $ tenemos que podemos encontrar $\displaystyle \tau, \rho > 0 $ suficientemente pequeños para que el cilindro
\[R_{\tau, \rho} := \left\{ \left(t,x\right)\in \R \times \R^{n} \; : \; \left|t-t_{0}\right|\leq \tau, \|x-x_{0}\|_{\infty}\leq \rho\right\} \subset D .\]
Es sencillo ver que como $\displaystyle f $ es continua y $\displaystyle R_{\tau, \rho} $ es un compacto, debe ser que $\displaystyle f $ está acotada en $\displaystyle R_{\tau, \rho} $ por lo que podemos definir
\[M := \max_{\left(t,x\right)\in R_{\tau, \rho}} \left|f\left(t,x\right)\right|>0 .\]
Además, como $\displaystyle f $ es localmente Lipschitz, podemos tomar $\displaystyle L $ como la constante de Lipschitz en el compacto $\displaystyle R_{\tau, \rho} $. Ahora, tomamos $\displaystyle \epsilon >0 $ tal que 
\[\epsilon < \min \left\{ \tau, \frac{\rho}{M}, \frac{1}{L}\right\}  ,\]
y definimos los conjuntos $\displaystyle E = \mathcal{C}\left(\left[t_{0}-\epsilon, t_{0} + \epsilon \right] , \R^{n}\right) $ y $\displaystyle X := \left\{ x \in E \; : \; \|x-x_{0}\|_{\infty } \leq \rho\right\}  $. Antes de aplicar el TPF tenemos que comprobar un par de cosas.
\begin{itemize}
\item En primer lugar, veamos que $\displaystyle T : X \to E $, es decir, $\displaystyle Ty $ es continua cuando $\displaystyle y \in X $. Tomamos $\displaystyle y \in X $ y $\displaystyle t_{1}, t_{2} \in I_{\epsilon } $ con $\displaystyle t_{1} > t_{2} $. Como $\displaystyle \epsilon < \tau $,
\[
\begin{split}
	\left\|\left(Ty\right)\left(t_{1}\right)-\left(Ty\right)\left(t_{2}\right)\right\| = & \left\|\int^{t_{1}}_{t_{0}} f\left(s, y\left(s\right)\right) \; d s - \int^{t_{2}}_{t_{0}} f\left(s, y\left(s\right)\right) \; d s\right\| = \left\|\int^{t_{1}}_{t_{2}} f\left(s, y\left(s\right)\right) \; d s\right\| \\
	\leq & \int^{t_{1}}_{t_{2}} \left\|f\underbrace{\left(s, y\left(s\right)\right)}_{\in R_{\tau, \rho}}\right\| \; d s \leq \int^{t_{1}}_{t_{2}} M \; d s = M \left|t_{1}-t_{2}\right|.
\end{split}
\]
	Así, tenemos que $\displaystyle T_{y} $ es continua. 
\item Veamos ahora que $\displaystyle T : X \to X $, es decir que $\displaystyle \|Ty-x_{0}\|_{\infty }\leq \rho $ para $\displaystyle y \in X $. Si $\displaystyle t \in I_{\epsilon } $, 
	\[
	\begin{split}
	\left\|\left(Ty\right)\left(t\right)-x_{0}\right\| = \left\|\int^{t}_{t_{0}} f\left(s, y\left(s\right)\right) \; d s\right\|\leq \int^{t}_{t_{0}} \left\|f\left(s, y\left(s\right)\right)\right\| \; d s \leq M \left|t-t_{0}\right| < M\epsilon <\rho .
	\end{split}
	\]
Por tanto, tenemos que $\displaystyle \|Ty-x_{0}\|_{\infty }\leq \rho $, por lo que $\displaystyle Ty \in X $. 
\item Veamos que $\displaystyle T $ es contractiva, es decir, que $\displaystyle \|Ty-Tz\|_{\infty }\leq K\|y-z\|_{\infty } $ con $\displaystyle K \in [0,1) $, $\displaystyle \forall y,z \in X$. Sean $\displaystyle y,z \in X $ y $\displaystyle t \in I_{\epsilon } $, entonces aplicando que $\displaystyle f $ es Lipschitz
	\[
	\begin{split}
		\left\|\left(Ty\right)\left(t\right)-\left(Tz\right)\left(t\right)\right\| = & \left\|\int^{t}_{t_{0}} f\left(s,y\left(s\right)\right)-f\left(s, z\left(s\right)\right) \; d s\right\|\leq \int^{t}_{t_{0}} \left\|f\left(s,y\left(s\right)\right)-f\left(s, z\left(s\right)\right)\right\| \; d s \\
		\leq & \int^{t}_{t_{0}} L \left\|y\left(s\right)-z\left(s\right)\right\| \; d s \leq \int^{t}_{t_{0}} L\max _{s\in I_{\epsilon }} \left\|y\left(s\right)-z\left(s\right)\right\| \; d s \\
		\leq & L \|y-z\|_{\infty } \left|t-t_{0}\right| < L \|y-z\|_{\infty }\epsilon .
	\end{split}
	\]
Así, tenemos que tomando máximos
\[\|Ty-Tz\|_{\infty }\leq L\epsilon \|y-z\|_{\infty } .\]
Como hemos tomado $\displaystyle \epsilon  $ tal que $\displaystyle L\epsilon < 1 $ tenemos que $\displaystyle T $ es contractiva. 
\end{itemize}
Por el TPF tenemos que existe una única $\displaystyle y \in X $ tal que $\displaystyle Ty = y $. Sin embargo, tenemos que ver que la unicidad de soluciones se verifica también en $\displaystyle E $, es decir, tenemos que ver que no hay soluciones extra en $\displaystyle E/X $. \\

Supongamos que existe una solución $\displaystyle z \in E/X $, es decir $\displaystyle \|z-x_{0}\|_{\infty} > \rho $. Así, existe $\displaystyle \xi \in I_{\epsilon } $ tal que $\displaystyle \left\|z\left(\xi\right)-x_{0}\left(\xi\right)\right\|>\rho $. Podemos suponer sin pérdida de generalidad que $\displaystyle \xi > t_{0} $. Definimos
\[\tilde{t} := \inf \left\{ t > t_{0} \; : \; \|z\left(\tilde{t}\right)-x_{0}\| > \rho\right\}\in (t_{0}, \xi]  .\]
Por continuidad tenemos que $\displaystyle \left|z\left(t\right)-x_{0}\right| \leq \rho $ en $\displaystyle \left[t_{0}, \tilde{t}\right]  $. Por tanto, $\displaystyle \left(t, z\left(t\right)\right)\in R $, $\displaystyle \forall t \in \left[t_{0}, \tilde{t}\right]  $. Deducimos que
\[
	\left\|z\left(\tilde{t}\right)-x_{0}\right\| =  \left\|\int^{\tilde{t}}_{t_{0}} f\left(s, z\left(s\right)\right) \; d s\right\| \leq \int^{\tilde{t}}_{t_{0}} \left\|f\left(s, z\left(s\right)\right)\right\| \; d s \leq M \left|\tilde{t}-t_{0}\right| < M\epsilon < \rho.
\]
Esto es una contradicción. 
\end{proof}
\begin{observation}
	\begin{enumerate}
	\item En la demostración necesitamos que $\displaystyle \epsilon < \max \left\{ \tau, \frac{\rho}{M}, \frac{1}{L}\right\}  $, por lo que $\displaystyle \epsilon  $ depende de $\displaystyle \left(t_{0}, x_{0}\right) $. 
	\item La dependencia de $\displaystyle \epsilon  $ respecto de la constante de Lipschitz no es necesaria. En efecto, se puede hacer con el corolario de que el TPF funciona con que una potencia del operador sea contractivo. Basta con ver que 
		\[\|T^{m}y-T^{m}z\|_{\infty } \leq \frac{\left(L\epsilon \right)^{m}}{m!}\|y-z\|_{\infty } .\]
		Como $\displaystyle \frac{\left(L\epsilon \right)^{m}}{m!}\to 0 $ si $\displaystyle m \to \infty $, se puede elegir $\displaystyle m \in \N $ (dependiendo de $\displaystyle \epsilon  $ y $\displaystyle L $) tal que $\displaystyle T^{m} $ es contracción. 
	\item Se puede elegir $\displaystyle \epsilon _{K} > 0 $ que sirva para toda condición inicial $\displaystyle \left(t_{0}, x_{0}\right)\in K $ con $\displaystyle K $ compacto. Basta con coger $\displaystyle \tilde{R} = \bigcup_{\left(t_{0}, x_{0}\right)\in K}R $ compacto y cogemos 
		\[0 < \epsilon < \max \left\{ \tau_{\left(t_{0}, x_{0}\right)}, \frac{\tau_{\left(t_{0}, x_{0}\right)}}{\max_{\tilde{R}}\|f\|}\right\} .\]
	\item No merece la pena optimizar el $\displaystyle \epsilon  $ del teorema eligiendo mejores cilindros. 
	\end{enumerate}
\end{observation}
\begin{eg}
Consideremos el PVI
\[
\begin{cases}
x' = 1 + x^{2} \\
x\left(0\right)=0
\end{cases}
.\]
Tenemos que en este caso $\displaystyle f\left(t,x\right)=1+x^{2} $ y $\displaystyle D = \R \times \R $. Buscamos el $\displaystyle \epsilon = \min \left\{ r, \frac{r}{\max_{R} \left|f\right|}\right\} =\min \left\{ r, \frac{r}{1+r^{2}}\right\} =\frac{1}{2} $ máximo que optimice el dominio de la solución. Tenemos que la solución explícita es $\displaystyle x\left(t\right)=\tan t $ para $\displaystyle t \in \left(-\frac{\pi }{2}, \frac{\pi }{2}\right) $ que es un intervalo mejor que el que habíamos conseguido.
\end{eg}

\begin{theorem}[Teorema Picard-Lindelöf global]
	Sea $\displaystyle f \in \mathcal{C}\left(\left[a,b\right] \times\R^{n}, \R^{n}\right) $ y global Lipschitz respecto a la segunda variable en $\displaystyle \left[a,b\right] \times \R^{n} $. Dado $\displaystyle \left(t_{0}, x_{0}\right) \in \left[a,b\right] \times \R^{n} $ y el PVI
	\[
	\begin{cases}
	x' = f\left(t,x\right)\\
	x\left(t_{0}\right)=x_{0}
	\end{cases}
	.\]
	Este sistema tiene solución única definida en todo $\displaystyle \left[a,b\right]  $. 	
\end{theorem}
\begin{proof}
	Basta con probar que el operador $\displaystyle T $ o una potencia suya es contractiva en $\displaystyle E = \mathcal{C}\left(\left[a,b\right] , \R^{n}\right) $. En la demostración anterior ya vimos que $\displaystyle T : E \to E $ está bien definida y que una potencia suya es contractiva. 
\end{proof}
\begin{colorary}
Sea $\displaystyle f \in \mathcal{C}\left(I \times \R^{n}, \R^{n}\right) $ globalmente Lipschitz respecto a la segunda variable en $\displaystyle I \times \R^{n} $, donde $\displaystyle I $ es un intervalo abierto. Dado $\displaystyle t_{0} \in I $ y $\displaystyle x_{0} \in \R^{n} $, el PVI tiene existencia y unicidad de solución definida en todo $\displaystyle I $.
\end{colorary}
\begin{proof}
	Consideramos $\displaystyle I =\left(a,b\right) $ un intervalo abierto con $\displaystyle -\infty\leq a < b \leq \infty $ y $\displaystyle t_{0} \in \left(a,b\right) $. Definimos una sucesión $\displaystyle \left\{ a_{k}\right\} _{k \in \N} $ estrictamente decreciente con $\displaystyle a_{k}\to a $. Análogamente, definimos una sucesión $\displaystyle \left\{ b_{k}\right\} _{k \in \N} $ estrictamente creciente con $\displaystyle b_{k} \to b $. 
	Así, obtenemos una sucesión de intervalos encajados cerrados tales que $\displaystyle t_{0} \in \left[a_{k}, b_{k}\right] \subset \left(a,b\right) $. En particular tenemos que $\displaystyle f \in \mathcal{C}\left(\left[a_{k}, b_{k}\right] \times \R^{n}, \R^{n}\right) $ es globalmente Lipschitz respecto a la segunda variable. 
	Aplicando el teorema Picard-Lindelöf global tenemos que existe una única solución al PVI, $\displaystyle x_{k} \in \mathcal{C}^{1}\left(\left[a_{k}, b_{k}\right] , \R^{n}\right) $. Es importante ver que por la unicidad de soluciones en $\displaystyle \left[a_{k}, b_{k}\right]  $ se tiene que
	\[x_{k+l}|_{\left[a_{k}, b_{k}\right] }=x_{k}, \; l \in \N .\]
	Dado $\displaystyle t\in I $, tenemos que existe $\displaystyle k \in \N $ tal que $\displaystyle t \in \left[a_{k}, b_{k}\right]  $ y entonces definimos $\displaystyle x\left(t\right):=x_{k}\left(t\right) $. Es sencillo comprobar que esta función es en efecto una solución. Veamos que es única. Supongamos que existe una función $\displaystyle y $ que también es solución. Entonces, tenemos que para un $\displaystyle k \in \N $, $\displaystyle y|_{\left[a_{k}, b_{k}\right] }=x $, por lo que debe ser que $\displaystyle y = x $. 
\end{proof}
\begin{observation}
	Realmente sólo necesitamos que $\displaystyle f $ sea globalmente Lipschitz en conjuntos de la forma $\displaystyle \left[a,b\right] \times \R^{n} $. Por tanto, no se requiere una constante de Lipschitz universal para todo el intervalo $\displaystyle I $.  
\end{observation}
\begin{eg}[Problema 2, Hoja 1]
Sea $\displaystyle f : \R \to \R $ y $\displaystyle -\infty = h_{0} < h_{1} < \cdots < h_{k} = \infty  $ tales que $\displaystyle f|_{\left(h_{j}, h_{j+1}\right)} $ es lineal. Si $\displaystyle f $ es continua, entonces es globalmente Lipschitz en $\displaystyle \R $. Podemos ver que la función $\displaystyle f\left(x\right)= \left|x\right| $ es un caso particular. 
Dado $\displaystyle x > y $ sean $\displaystyle i\geq j $ tales que $\displaystyle x\in \left[h_{i}, h_{i+1}\right]  $ e $\displaystyle y \in \left[h_{j}, h_{j+1}\right]  $. Así, tenemos que
\[
\begin{split}
	\left|f\left(x\right)-f\left(y\right)\right| = & \left|f\left(x\right)-f\left(h_{i}\right)+f\left(h_{i}\right)-f\left(h_{i-1}\right)+f\left(h_{i-1}\right)-\cdots +f\left(h_{j+1}\right)-f\left(y\right)\right|\\
	\leq & \left|f\left(x\right)-f\left(h_{1}\right)\right| + \cdots + \left|f\left(h_{j+1}\right)-f\left(y\right)\right| \\
	= & \left|a_{i}\left(x-h_{i}\right)\right|+ \cdots + \left|a_{j}\left(h_{j+1}\right)-y\right|= \left|a_{i}\right|\left(x-h_{i}\right)+ \cdots + \left|a_{j}\right|\left(h_{j+1}-y\right)\\
	\leq & \max_{l \in \left\{ j, \ldots, i\right\} } \left|a_{l}\right| \left|x-y\right|.
\end{split}
\]
Así, para terminar cogemos $\displaystyle K = \max_{l \in \left\{ 0, \ldots, k-1\right\} } \left|a_{l}\right| $ como constante de Lipschitz. 
\end{eg}
\begin{theorem}[Teorema de Peano]
	Dado el PVI
	\[
	\begin{cases}
	x'= f\left(t,x\right)\\
	x\left(t_{0}\right)=x_{0}
	\end{cases}
	,\]
	con $\displaystyle f \in \mathcal{C}\left(D, \R^{n}\right) $ en un abierto $\displaystyle D \subset \R\times \R^{n} $ y $\displaystyle \left(t_{0}, x_{0}\right) \in D $, existe una solución $\displaystyle x \in \mathcal{C}^{1}\left(\left[t_{0}-\epsilon, t_{0} + \epsilon \right] , \R^{n}\right) $ del PVI para cierto $\displaystyle \epsilon > 0 $. 
\end{theorem}
\begin{proof}
No lo vamos a demostrar porque es bastante laborioso pero damos tres estrategias para demostrarlo utilizando el teorema de Ascoli-Arzelá. 
\begin{enumerate}
	\item Aproximar $\displaystyle f $ por polinomios (que son localmente Lipschitz) y resolvemos el sistema $\displaystyle x' = f_{\epsilon }\left(t,x\right) $ obteniendo una solución $\displaystyle x_{\epsilon } $ por el teorema de Picard-Lindelöf. Tendríamos que probar que $\displaystyle x_{\epsilon } \xrightarrow{\epsilon \to 0} x$. 
	\item Construir las poligonales de Euler. 
	\item Utilizar un teorema del punto fijo (uno mejor que el que tenemos). Concretamente, usar el teorema del punto fijo de Schauder.
\end{enumerate}
\end{proof}

